\section{Circuit Complexity --- Fan Out}
%%% generated by the blueprint pipeline from TCSlib.Complexity.CircuitComplexity.FanOut
\subsection*{Circuit Complexity --- Fan Out: Introduction}

Declarations of the Lean module \texttt{TCSlib.Complexity.CircuitComplexity.FanOut}.

\subsection*{Circuit Complexity --- Fan Out: Results}

\begin{definition}[Fanout of a vertex in a gate list]
\lean{BoolCircuit.fanoutIn}
\label{BoolCircuit.fanoutIn}
\leanok
\uses{BoolCircuit.DAGGate}
For a list of directed acyclic graph gates and a natural-number vertex $v$, define the
fanout of $v$ to be the number of gates in the list whose argument list contains $v$.
Each such gate is counted once, even if $v$ occurs more than once among its arguments.
\end{definition}

\begin{definition}[Fanout of a circuit vertex]
\lean{BoolCircuit.DAGCircuit.fanout}
\label{BoolCircuit.DAGCircuit.fanout}
\leanok
\uses{BoolCircuit.DAGCircuit, BoolCircuit.fanoutIn}
For a directed acyclic circuit and a natural-number vertex $v$, define its fanout to be
the number of gates in the circuit whose argument list contains $v$.
\end{definition}

\begin{definition}[Circuit with fanout at most two]
\lean{BoolCircuit.DAGCircuit.HasFanoutTwo}
\label{BoolCircuit.DAGCircuit.HasFanoutTwo}
\leanok
\uses{BoolCircuit.DAGCircuit, BoolCircuit.DAGCircuit.fanout}
A directed acyclic circuit has fanout at most two if every natural-number vertex has
fanout at most $2$.
\end{definition}

\begin{theorem}[Fan-out under concatenation]
\lean{BoolCircuit.fanoutIn_append}
\label{BoolCircuit.fanoutIn_append}
\leanok
\uses{BoolCircuit.DAGGate, BoolCircuit.fanoutIn}
\difficulty{1}
For any two finite lists of gates $G$ and $E$ and any vertex $v$, the number of gates in
$G$ followed by $E$ that read $v$ equals the number in $G$ that read $v$ plus the number
in $E$ that read $v$.
\end{theorem}

\begin{theorem}[Fanout in a single gate]
\lean{BoolCircuit.fanoutIn_singleton}
\label{BoolCircuit.fanoutIn_singleton}
\leanok
\uses{BoolCircuit.DAGGate, BoolCircuit.fanoutIn}
\difficulty{4}
For any Boolean gate $g$ and vertex $v$, the fanout of $v$ among the one-element list
containing $g$ is $1$ if $g$ reads $v$, and $0$ otherwise.
\end{theorem}

\begin{theorem}[Zero fan-out beyond the circuit]
\lean{BoolCircuit.fanoutIn_eq_zero_of_le}
\label{BoolCircuit.fanoutIn_eq_zero_of_le}
\leanok
\uses{BoolCircuit.DAGGate, BoolCircuit.GatesAcyclic, BoolCircuit.fanoutIn}
\difficulty{4}
Let $n$ be a natural number and let $G$ be a finite acyclic list of gates on $n$ input
vertices, with each gate reading only vertices preceding it. If $n+|G|\le v$, then no
gate in $G$ reads vertex $v$; equivalently, the fan-out of $v$ within $G$ is zero.
\end{theorem}

\begin{theorem}[Acyclicity after appending gates]
\lean{BoolCircuit.gatesAcyclic_append_of_lt}
\label{BoolCircuit.gatesAcyclic_append_of_lt}
\leanok
\uses{BoolCircuit.DAGGate, BoolCircuit.GatesAcyclic}
\difficulty{4}
Let $n$ be a natural number and let $gs$ be an acyclic finite list of gates starting
from index bound $n$. If every input index read by every gate in a finite extension
$ext$ is less than $n+|gs|$, then the concatenation $gs$ followed by $ext$ is acyclic
starting from $n$.
\end{theorem}

\begin{definition}[Single gate fan-out expansion step]
\lean{BoolCircuit.foStep}
\label{BoolCircuit.foStep}
\leanok
\uses{BoolCircuit.DAGGate}
Given a number $n$, a state consisting of a gate list, a map from natural numbers to
natural numbers, and a counter, and a gate $g$, write the state as $(G,t,N)$. The step
appends to $G$ the gate with the same kind as $g$ and with each argument $a$ replaced by
$t(a)$, followed by one AND gate with sole argument $t(a)$ for each argument $a$ of $g$.
It replaces $t(v)$ by $n+|G|$ when $v=N$, by $n+|G|+1+i$ when $v$ occurs in the argument
list of $g$ at index $i$, and otherwise leaves $t(v)$ unchanged; the new counter is
$N+1$.
\end{definition}

\begin{definition}[Iterated fan-out expansion]
\lean{BoolCircuit.foGates}
\label{BoolCircuit.foGates}
\leanok
\uses{BoolCircuit.DAGGate, BoolCircuit.foStep}
Given a natural number $n$ and a list of gates, start with the empty gate list, the
identity map on natural numbers, and counter $n$, then process the gates in order by
applying the single-gate fan-out expansion step with parameter $n$. The resulting state
is the iterated fan-out expansion of the list.
\end{definition}

\begin{definition}[Fan-out expansion invariant]
\lean{BoolCircuit.FoInv}
\label{BoolCircuit.FoInv}
\leanok
\uses{BoolCircuit.DAGGate, BoolCircuit.DAGGate.eval, BoolCircuit.GatesAcyclic, BoolCircuit.fanoutIn, BoolCircuit.runWith, BoolCircuit.vertexValue}
For a natural number $n$, an old list of gates, and a state $(G,t,N)$, the fan-out
expansion invariant consists of the following properties. The counter is $n$ plus the
length of the old list; $G$ is acyclic; for every $v<n+|	ext{old}|$, its image $t(v)$ is
less than $n+|G|$, and $t$ is injective on those indices. For every Boolean input $x$ of
length $n$ and every such $v$, the value at vertex $t(v)$ in $G$ equals the value of
vertex $v$ when the old gates are run on $x$. Every vertex has fan-out at most $2$, and
each $t(v)$ in that range has fan-out $0$. Every gate in $G$ has distinct arguments, a
NOT gate has exactly one argument, and every gate has at most two arguments. Finally,
$|G|\le 3|\text{old}|$.
\end{definition}

\begin{theorem}[Invariants for the empty circuit]
\lean{BoolCircuit.foInv_nil}
\label{BoolCircuit.foInv_nil}
\leanok
\uses{BoolCircuit.FoInv, BoolCircuit.GatesAcyclic.nil, BoolCircuit.fanoutIn}
\difficulty{3}
For every nonnegative integer $n$, the empty original gate list together with the empty
new gate list, the identity vertex map, and count $n$ satisfies the fanout reduction
invariants: the count is $n$; the new list is acyclic; the map is injective and sends
each vertex below $n$ below $n$; computed values agree for every Boolean input; all
fanouts are at most $2$ and mapped vertices have fanout zero; the gate conditions hold;
and the new list has length at most zero.
\end{theorem}

\begin{theorem}[Preservation of fanout reduction invariants]
\lean{BoolCircuit.foInv_step}
\label{BoolCircuit.foInv_step}
\leanok
\uses{BoolCircuit.DAGGate, BoolCircuit.DAGGate.eval, BoolCircuit.DAGGate.eval_remap, BoolCircuit.DAGGate.eval_singleton, BoolCircuit.FoInv, BoolCircuit.GatesAcyclic, BoolCircuit.fanoutIn, BoolCircuit.fanoutIn_append, BoolCircuit.fanoutIn_eq_zero_of_le, BoolCircuit.foStep, BoolCircuit.gatesAcyclic_append_of_lt, BoolCircuit.runWith, BoolCircuit.runWith_append, BoolCircuit.runWith_getD_gate, BoolCircuit.runWith_getD_last, BoolCircuit.runWith_getD_of_lt, BoolCircuit.runWith_getD_take, BoolCircuit.vertexValue, BoolCircuit.vertexValue_append, BoolCircuit.vertexValue_last}
\difficulty{7}
Let $n$ be a nonnegative, integer, let $old$ be a finite list of Boolean gates, and let
$g$ be another Boolean gate. Suppose a gate list, vertex map, and count satisfy the
fanout reduction invariants for $old$. Suppose also that appending $g$ to $old$
preserves acyclicity, and that $g$ has distinct arguments, has at most two arguments,
and has exactly one argument if it is a NOT gate. Then applying one fanout reduction
step with $g$ yields a list, map, and count satisfying the fanout reduction invariants
for the list obtained by appending $g$ to $old$.
\end{theorem}

\begin{theorem}[Invariants of the fanout reduction]
\lean{BoolCircuit.foGates_inv}
\label{BoolCircuit.foGates_inv}
\leanok
\uses{BoolCircuit.DAGGate, BoolCircuit.FoInv, BoolCircuit.GatesAcyclic, BoolCircuit.foGates, BoolCircuit.foInv_nil, BoolCircuit.foInv_step}
\difficulty{5}
Let $n$ be a nonnegative integer and let $old$ be a finite list of Boolean gates.
Suppose $old$ is acyclic, meaning every gate at position $i$ reads only vertices with
index less than $n+i$, and every gate has distinct arguments, a NOT gate has exactly one
argument, and every gate has at most two arguments. The fanout reduction of $old$
produces a list of gates, a map assigning a new vertex to each old vertex, and a count.
The count equals $n+|old|$; the new list is acyclic; every old vertex below $n+|old|$
maps below $n$ plus the new list’s length, and the map is injective on those vertices.
For every Boolean input on $n$ variables, the value at each mapped vertex equals the
corresponding value in the original computation (using false for indices beyond the
computed list). Every vertex has fanout at most $2$, each mapped old vertex has fanout
zero, every new gate has distinct arguments and at most two arguments, and each NOT gate
has exactly one argument. The new list has length at most $3|old|$.
\end{theorem}

\begin{theorem}[Gate properties of fan-in-two circuits]
\lean{BoolCircuit.gates_ok_of_isFaninTwo}
\label{BoolCircuit.gates_ok_of_isFaninTwo}
\leanok
\uses{BoolCircuit.DAGCircuit, BoolCircuit.DAGCircuit.IsFaninTwo}
\difficulty{1}
Let $C$ be a finite directed acyclic circuit. Suppose every gate of $C$ has pairwise
distinct input vertices, every NOT gate has exactly one input, and every gate has at
most two inputs. Then each gate has pairwise distinct inputs, a NOT gate has exactly one
input, and the gate has at most two inputs.
\end{theorem}

\begin{definition}[Circuit with fan-out at most two]
\lean{BoolCircuit.DAGCircuit.fanoutTwo}
\label{BoolCircuit.DAGCircuit.fanoutTwo}
\leanok
\uses{BoolCircuit.DAGCircuit, BoolCircuit.DAGCircuit.IsFaninTwo, BoolCircuit.foGates, BoolCircuit.foGates_inv, BoolCircuit.gates_ok_of_isFaninTwo}
Let $C$ be a finite acyclic Boolean circuit with $n$ input vertices, whose gates have
distinct arguments, whose NOT gates each have exactly one argument, and whose gates each
have at most two arguments. The circuit obtained by iterated fan-out expansion of its
gate list has the expanded gate list and the output vertex obtained by applying the
expansion's map to the original output.
\end{definition}

\begin{theorem}[Evaluation preserved under fanout reduction]
\lean{BoolCircuit.DAGCircuit.fanoutTwo_eval}
\label{BoolCircuit.DAGCircuit.fanoutTwo_eval}
\leanok
\uses{BoolCircuit.DAGCircuit.eval, BoolCircuit.DAGCircuit.fanoutTwo, BoolCircuit.foGates_inv, BoolCircuit.gates_ok_of_isFaninTwo}
\difficulty{1}
Let $C$ be a Boolean circuit with $n$ input vertices, whose gates are in topological
order and whose designated output is a circuit vertex. Suppose each gate reads distinct
vertices, each negation gate reads exactly one vertex, and every gate reads at most two
vertices. Then, for every assignment $x$ of Boolean values to the input vertices, the
circuit obtained from $C$ by the fanout-two construction has the same output value as
$C$.
\end{theorem}

\begin{theorem}[Fanin bound after fanout reduction]
\lean{BoolCircuit.DAGCircuit.fanoutTwo_isFaninTwo}
\label{BoolCircuit.DAGCircuit.fanoutTwo_isFaninTwo}
\leanok
\uses{BoolCircuit.DAGCircuit.IsFaninTwo, BoolCircuit.DAGCircuit.fanoutTwo, BoolCircuit.foGates_inv, BoolCircuit.gates_ok_of_isFaninTwo}
\difficulty{3}
Let $C$ be a finite directed acyclic Boolean circuit in which every gate has distinct
arguments, every NOT gate has exactly one argument, and every gate has at most two
arguments. The circuit obtained by the fanout reduction construction also has distinct
arguments at every gate, has exactly one argument at each NOT gate, and has at most two
arguments at every gate.
\end{theorem}

\begin{theorem}[Fanout bound after fanout reduction]
\lean{BoolCircuit.DAGCircuit.fanoutTwo_hasFanoutTwo}
\label{BoolCircuit.DAGCircuit.fanoutTwo_hasFanoutTwo}
\leanok
\uses{BoolCircuit.DAGCircuit.HasFanoutTwo, BoolCircuit.DAGCircuit.fanoutTwo, BoolCircuit.foGates_inv, BoolCircuit.gates_ok_of_isFaninTwo}
\difficulty{1}
Let $C$ be a finite directed acyclic Boolean circuit in which every gate has distinct
arguments, every NOT gate has exactly one argument, and every gate has at most two
arguments. In the circuit obtained by the fanout reduction construction, every vertex is
read by at most two gates.
\end{theorem}

\begin{theorem}[Size bound after fanout reduction]
\lean{BoolCircuit.DAGCircuit.fanoutTwo_size_le}
\label{BoolCircuit.DAGCircuit.fanoutTwo_size_le}
\leanok
\uses{BoolCircuit.DAGCircuit.fanoutTwo, BoolCircuit.DAGCircuit.size, BoolCircuit.foGates, BoolCircuit.foGates_inv, BoolCircuit.gates_ok_of_isFaninTwo}
\difficulty{3}
Let $C$ be a finite directed acyclic Boolean circuit in which every gate has distinct
arguments, every NOT gate has exactly one argument, and every gate has at most two
arguments. The size of the circuit obtained by the fanout reduction construction is at
most $n+3m$, where $n$ is the number of input vertices and $m$ is the number of gates of
$C$.
\end{theorem}

\begin{theorem}[Bounded-fanout equivalent circuit]
\lean{BoolCircuit.DAGCircuit.exists_fanoutTwo}
\label{BoolCircuit.DAGCircuit.exists_fanoutTwo}
\leanok
\uses{BoolCircuit.DAGCircuit, BoolCircuit.DAGCircuit.HasFanoutTwo, BoolCircuit.DAGCircuit.IsFaninTwo, BoolCircuit.DAGCircuit.eval, BoolCircuit.DAGCircuit.fanoutTwo, BoolCircuit.DAGCircuit.fanoutTwo_eval, BoolCircuit.DAGCircuit.fanoutTwo_hasFanoutTwo, BoolCircuit.DAGCircuit.fanoutTwo_isFaninTwo, BoolCircuit.DAGCircuit.fanoutTwo_size_le, BoolCircuit.DAGCircuit.size}
\difficulty{3}
Let $C$ be a directed acyclic Boolean circuit on $n$ inputs in which each gate has no
repeated arguments, every NOT gate has exactly one argument, and every gate has at most
two arguments. There is a circuit $C'$ on the same inputs satisfying these conditions,
such that each vertex is an argument of at most two gates, $C'$ computes the same
Boolean function as $C$, and the size of $C'$ is at most three times the size of $C$.
\end{theorem}
